\documentclass[10pt]{article}
\usepackage[margin=1.6cm]{geometry}
\usepackage{amsmath,amssymb,xcolor,booktabs,longtable,array,parskip,microtype}
\definecolor{chg}{HTML}{C0392B}
\setcounter{secnumdepth}{1}
\begin{document}
\begin{center}{\LARGE\bfseries A Full Per-Token Trace of \texttt{\_s2\_closer}}\\[3pt]
{\large every s-op as a row, every token as a column ---}\\[1pt]
{\large\ttfamily HI = the boy that chases the dog swims}\\[3pt]
{\small Generated by \texttt{\_gen\_closer\_handout.py}; all values asserted equal to the module and \texttt{grammar\_oracle.py}.}\end{center}
\hrule\vspace{4pt}
\section{Stage 1 --- encoding and part of speech}
Each column is a token \emph{position} (note \emph{the} appears at both 0 and 4). Every row is one s-op: a length-7 array holding one value per position. Lexical-class codes: 0=D, 1=N\_sing, 2=N\_prop, 3=Adj, 4=Adv, 5=THAT, 6=V, 7=OTHER. Verb-role codes: 0=non-verb, 1=intransitive, 2=+object DP, 3=+clause, 4=object-gap.
\begin{center}\footnotesize
\setlength{\tabcolsep}{4pt}\renewcommand{\arraystretch}{1.15}
\begin{longtable}{@{\extracolsep{\fill}}l|ccccccc@{}}
\toprule \textbf{s-op} & the$_{0}$ & boy$_{1}$ & that$_{2}$ & chases$_{3}$ & the$_{4}$ & dog$_{5}$ & swims$_{6}$ \\ \midrule \endfirsthead
\toprule \textbf{s-op} & the$_{0}$ & boy$_{1}$ & that$_{2}$ & chases$_{3}$ & the$_{4}$ & dog$_{5}$ & swims$_{6}$ \\ \midrule \endhead \bottomrule \endfoot
\multicolumn{8}{@{}p{0.98\textwidth}@{}}{\itshape The fixed word$\to$id lookup $\texttt{\_W2I}$ (a token embedding). Each surface word becomes its unique integer id; repeated words (\emph{the}) get the \emph{same} id.} \\[1pt]
\texttt{hi\_ids} & 0 & 5 & 53 & 33 & 0 & 3 & 45 \\
\addlinespace[3pt]
\multicolumn{8}{@{}p{0.98\textwidth}@{}}{\itshape The position s-op $[0,1,\dots,n{-}1]$ from $\texttt{indices}$. The whole program's output is $\texttt{idx}+\text{disp}$, so this is the identity we perturb.} \\[1pt]
\texttt{idx} & 0 & 1 & 2 & 3 & 4 & 5 & 6 \\
\addlinespace[3pt]
\multicolumn{8}{@{}p{0.98\textwidth}@{}}{\itshape A second embedding $\texttt{tok\_map(hi\_ids,\,\_lexclass)}$: id$\to$coarse POS class. Shown by name; the underlying values are the integer codes above.} \\[1pt]
\texttt{lc} & D & N & THAT & V & D & N & V \\
\addlinespace[3pt]
\multicolumn{8}{@{}p{0.98\textwidth}@{}}{\itshape Pure lexical flag: 1 iff the word is in $\texttt{V\_intrans}$. Only \emph{swims} qualifies. Used later to tell an intransitive verb from an object-gap verb.} \\[1pt]
\texttt{is\_intr} & 0 & 0 & 0 & 0 & 0 & 0 & 1 \\
\addlinespace[3pt]
\multicolumn{8}{@{}p{0.98\textwidth}@{}}{\itshape Backward neighbour read $\texttt{\_read\_off(lc,-1)}=lc[i{-}1]$ (one attention head keyed on relative offset). $-1$ = fell off the left edge.} \\[1pt]
\texttt{prev\_lc} & -- & D & N & THAT & V & D & N \\
\addlinespace[3pt]
\multicolumn{8}{@{}p{0.98\textwidth}@{}}{\itshape $lc[i{-}2]$. Feeds the `adjective-core' relative-clause test far below ($\texttt{has\_adj\_core}$).} \\[1pt]
\texttt{prev2\_lc} & -- & -- & D & N & THAT & V & D \\
\addlinespace[3pt]
\multicolumn{8}{@{}p{0.98\textwidth}@{}}{\itshape Forward read $lc[i{+}1]$. This is the read that relies on the `HI block is in the past' convention (it looks at a later position, legal only because the whole HI block precedes the output region).} \\[1pt]
\texttt{next\_lc} & N & THAT & V & D & N & V & -- \\
\addlinespace[3pt]
\multicolumn{8}{@{}p{0.98\textwidth}@{}}{\itshape $lc[i{+}2]$, used to skip one optional post-verb adverb.} \\[1pt]
\texttt{next2\_lc} & THAT & V & D & N & V & -- & -- \\
\addlinespace[3pt]
\multicolumn{8}{@{}p{0.98\textwidth}@{}}{\itshape $(lc=\text{THAT})$. Marks the clause \textbf{openers} --- the only visible bracket-starts in the string. Fires only at position 2.} \\[1pt]
\texttt{isthat} & 0 & 0 & 1 & 0 & 0 & 0 & 0 \\
\addlinespace[3pt]
\multicolumn{8}{@{}p{0.98\textwidth}@{}}{\itshape $\texttt{isthat}\ \&\ (prev\_lc=\text{N})$: a \emph{that} right after a common noun is a \textbf{relative} complementizer. Here \emph{boy}(1) precedes \emph{that}(2), so this fires.} \\[1pt]
\texttt{is\_crel} & 0 & 0 & 1 & 0 & 0 & 0 & 0 \\
\addlinespace[3pt]
\multicolumn{8}{@{}p{0.98\textwidth}@{}}{\itshape $\texttt{isthat}\ \&\ \lnot\texttt{is\_crel}$: the \emph{sentential} \emph{that}. Empty here --- our only \emph{that} is relative.} \\[1pt]
\texttt{is\_c} & 0 & 0 & 0 & 0 & 0 & 0 & 0 \\
\addlinespace[3pt]
\multicolumn{8}{@{}p{0.98\textwidth}@{}}{\itshape $(lc=\text{V})$. True at \emph{chases}(3) and \emph{swims}(6).} \\[1pt]
\texttt{isverb} & 0 & 0 & 0 & 1 & 0 & 0 & 1 \\
\addlinespace[3pt]
\multicolumn{8}{@{}p{0.98\textwidth}@{}}{\itshape $(lc=\text{Adv})$. Empty here (no adverb).} \\[1pt]
\texttt{isadv} & 0 & 0 & 0 & 0 & 0 & 0 & 0 \\
\addlinespace[3pt]
\multicolumn{8}{@{}p{0.98\textwidth}@{}}{\itshape `next non-adverb class': $\texttt{where(next\_lc=Adv,\,next2\_lc,\,next\_lc)}$. Lets a verb see its complement's first token even across an adverb. With no adverb it equals $\texttt{next\_lc}$.} \\[1pt]
\texttt{nn\_lc} & N & THAT & V & D & N & V & -- \\
\addlinespace[3pt]
\multicolumn{8}{@{}p{0.98\textwidth}@{}}{\itshape The verb role, from $\texttt{nn\_lc}$: \emph{chases}(3) sees \emph{the}=D $\Rightarrow$ 2 (transitive); \emph{swims}(6) is lexically intransitive $\Rightarrow$ 1. Roles 2/3 head a flip VP; 1/4 do not.} \\[1pt]
\texttt{role} & 0 & 0 & 0 & 2 & 0 & 0 & 1 \\
\addlinespace[3pt]
\end{longtable}\end{center}
\section{Stage 2 scaffolding --- one-shot local span pieces}
These are computed once, \emph{before} the loop; each is a fixed-offset lookahead, no iteration.
\begin{center}\footnotesize
\setlength{\tabcolsep}{4pt}\renewcommand{\arraystretch}{1.15}
\begin{longtable}{@{\extracolsep{\fill}}l|ccccccc@{}}
\toprule \textbf{s-op} & the$_{0}$ & boy$_{1}$ & that$_{2}$ & chases$_{3}$ & the$_{4}$ & dog$_{5}$ & swims$_{6}$ \\ \midrule \endfirsthead
\toprule \textbf{s-op} & the$_{0}$ & boy$_{1}$ & that$_{2}$ & chases$_{3}$ & the$_{4}$ & dog$_{5}$ & swims$_{6}$ \\ \midrule \endhead \bottomrule \endfoot
\multicolumn{8}{@{}p{0.98\textwidth}@{}}{\itshape Verb head-block end: $i{+}1$ for a verb\,+\,adverb, else $i$. No adverbs here, so it equals $\texttt{idx}$. This is the `head' side of a VP flip.} \\[1pt]
\texttt{hb\_end} & 0 & 1 & 2 & 3 & 4 & 5 & 6 \\
\addlinespace[3pt]
\multicolumn{8}{@{}p{0.98\textwidth}@{}}{\itshape Where a DP ends \emph{ignoring} relative clauses: $D\,(Adj)\,N$. \emph{the}(0)$\to$\emph{boy}(1); \emph{the}(4)$\to$\emph{dog}(5). Note it does \emph{not} yet know about the relative clause on \emph{boy}.} \\[1pt]
\texttt{dp\_end\_flat} & 1 & 1 & 2 & 3 & 5 & 5 & 6 \\
\addlinespace[3pt]
\multicolumn{8}{@{}p{0.98\textwidth}@{}}{\itshape $\texttt{dp\_end\_flat}+1$: the slot just past a DP's noun, where a relative \emph{that} would sit. For \emph{the}(0) this is 2.} \\[1pt]
\texttt{rc\_that\_pos} & 2 & 2 & 3 & 4 & 6 & 6 & 7 \\
\addlinespace[3pt]
\multicolumn{8}{@{}p{0.98\textwidth}@{}}{\itshape $\texttt{index\_select(lc,\,rc\_that\_pos)}$: the class actually sitting at that slot (one cross-position lookahead). At position 0 it reads $lc[2]=\text{THAT}$.} \\[1pt]
\texttt{lc\_at\_rc} & THAT & THAT & V & D & V & V & -- \\
\addlinespace[3pt]
\multicolumn{8}{@{}p{0.98\textwidth}@{}}{\itshape $(lc=D)\ \&\ (\texttt{lc\_at\_rc}=\text{THAT})$: this DP carries a relative clause. True at \emph{the}(0) --- `the boy \textbf{that}\,\dots'.} \\[1pt]
\texttt{dp\_has\_rc} & 1 & 0 & 0 & 0 & 0 & 0 & 0 \\
\addlinespace[3pt]
\end{longtable}\end{center}
\section{Stage 2 relaxation --- the loop, all 10 rounds}
Four END pointers are seeded to a lower bound (`each ends at itself') and refined by 10 sweeps. A \textcolor{chg}{\textbf{red}} cell changed from the row above; once a table stops turning red it has reached its fixed point. Only the openers'/verbs' columns carry meaning --- other columns hold the harmless $\texttt{idx}$ filler.
\paragraph{\texttt{dpfull}.} DP end including a trailing relative clause. Watch column \emph{the}$_0$: at $\texttt{init}$ it is the flat 1; it cannot finish until the relative clause's end is known, so it lags one sweep behind $\texttt{rce}$ and only settles at \textbf{R2}.
\begin{center}\footnotesize\setlength{\tabcolsep}{4pt}\renewcommand{\arraystretch}{1.1}
\begin{tabular}{@{}l|ccccccc@{}}
\toprule \textbf{round} & the$_{0}$ & boy$_{1}$ & that$_{2}$ & chases$_{3}$ & the$_{4}$ & dog$_{5}$ & swims$_{6}$ \\ \midrule
init & 1 & 1 & 2 & 3 & 5 & 5 & 6 \\
R1 & \textcolor{chg}{\textbf{2}} & 1 & 2 & 3 & 5 & 5 & 6 \\
R2 & \textcolor{chg}{\textbf{5}} & 1 & 2 & 3 & 5 & 5 & 6 \\
R3 & 5 & 1 & 2 & 3 & 5 & 5 & 6 \\
R4 & 5 & 1 & 2 & 3 & 5 & 5 & 6 \\
R5 & 5 & 1 & 2 & 3 & 5 & 5 & 6 \\
R6 & 5 & 1 & 2 & 3 & 5 & 5 & 6 \\
R7 & 5 & 1 & 2 & 3 & 5 & 5 & 6 \\
R8 & 5 & 1 & 2 & 3 & 5 & 5 & 6 \\
R9 & 5 & 1 & 2 & 3 & 5 & 5 & 6 \\
R10 & 5 & 1 & 2 & 3 & 5 & 5 & 6 \\
\bottomrule\end{tabular}\end{center}
\paragraph{\texttt{vpe}.} VP end (verb\,+\,complement). \emph{chases}$_3$ jumps to its object DP's end (5) in R1; \emph{swims}$_6$ is intransitive so it never moves off 6.
\begin{center}\footnotesize\setlength{\tabcolsep}{4pt}\renewcommand{\arraystretch}{1.1}
\begin{tabular}{@{}l|ccccccc@{}}
\toprule \textbf{round} & the$_{0}$ & boy$_{1}$ & that$_{2}$ & chases$_{3}$ & the$_{4}$ & dog$_{5}$ & swims$_{6}$ \\ \midrule
init & 0 & 1 & 2 & 3 & 4 & 5 & 6 \\
R1 & 0 & 1 & 2 & \textcolor{chg}{\textbf{5}} & 4 & 5 & 6 \\
R2 & 0 & 1 & 2 & 5 & 4 & 5 & 6 \\
R3 & 0 & 1 & 2 & 5 & 4 & 5 & 6 \\
R4 & 0 & 1 & 2 & 5 & 4 & 5 & 6 \\
R5 & 0 & 1 & 2 & 5 & 4 & 5 & 6 \\
R6 & 0 & 1 & 2 & 5 & 4 & 5 & 6 \\
R7 & 0 & 1 & 2 & 5 & 4 & 5 & 6 \\
R8 & 0 & 1 & 2 & 5 & 4 & 5 & 6 \\
R9 & 0 & 1 & 2 & 5 & 4 & 5 & 6 \\
R10 & 0 & 1 & 2 & 5 & 4 & 5 & 6 \\
\bottomrule\end{tabular}\end{center}
\paragraph{\texttt{rce}.} $\texttt{CP\_rel}$ end. \emph{that}$_2$ is a subject-gap relative, so its end = the VP right after it ($\texttt{vpe}$ at position 3), reaching 5 in R1.
\begin{center}\footnotesize\setlength{\tabcolsep}{4pt}\renewcommand{\arraystretch}{1.1}
\begin{tabular}{@{}l|ccccccc@{}}
\toprule \textbf{round} & the$_{0}$ & boy$_{1}$ & that$_{2}$ & chases$_{3}$ & the$_{4}$ & dog$_{5}$ & swims$_{6}$ \\ \midrule
init & 0 & 1 & 2 & 3 & 4 & 5 & 6 \\
R1 & 0 & 1 & \textcolor{chg}{\textbf{5}} & 3 & 4 & 5 & 6 \\
R2 & 0 & 1 & 5 & 3 & 4 & 5 & 6 \\
R3 & 0 & 1 & 5 & 3 & 4 & 5 & 6 \\
R4 & 0 & 1 & 5 & 3 & 4 & 5 & 6 \\
R5 & 0 & 1 & 5 & 3 & 4 & 5 & 6 \\
R6 & 0 & 1 & 5 & 3 & 4 & 5 & 6 \\
R7 & 0 & 1 & 5 & 3 & 4 & 5 & 6 \\
R8 & 0 & 1 & 5 & 3 & 4 & 5 & 6 \\
R9 & 0 & 1 & 5 & 3 & 4 & 5 & 6 \\
R10 & 0 & 1 & 5 & 3 & 4 & 5 & 6 \\
\bottomrule\end{tabular}\end{center}
\paragraph{\texttt{cpe}.} $\texttt{CP\_sent}$ end. Never active here (no sentential \emph{that}), so it stays at the $\texttt{idx}$ seed every round --- a useful negative example.
\begin{center}\footnotesize\setlength{\tabcolsep}{4pt}\renewcommand{\arraystretch}{1.1}
\begin{tabular}{@{}l|ccccccc@{}}
\toprule \textbf{round} & the$_{0}$ & boy$_{1}$ & that$_{2}$ & chases$_{3}$ & the$_{4}$ & dog$_{5}$ & swims$_{6}$ \\ \midrule
init & 0 & 1 & 2 & 3 & 4 & 5 & 6 \\
R1 & 0 & 1 & 2 & 3 & 4 & 5 & 6 \\
R2 & 0 & 1 & 2 & 3 & 4 & 5 & 6 \\
R3 & 0 & 1 & 2 & 3 & 4 & 5 & 6 \\
R4 & 0 & 1 & 2 & 3 & 4 & 5 & 6 \\
R5 & 0 & 1 & 2 & 3 & 4 & 5 & 6 \\
R6 & 0 & 1 & 2 & 3 & 4 & 5 & 6 \\
R7 & 0 & 1 & 2 & 3 & 4 & 5 & 6 \\
R8 & 0 & 1 & 2 & 3 & 4 & 5 & 6 \\
R9 & 0 & 1 & 2 & 3 & 4 & 5 & 6 \\
R10 & 0 & 1 & 2 & 3 & 4 & 5 & 6 \\
\bottomrule\end{tabular}\end{center}
After R2 nothing changes: the depth-1 sentence has reached the fixed point with rounds 3--10 merely confirming it. A depth-$d$ sentence would keep turning red through about round $d{+}1$.
\section{Stage 2 output --- clause ends and depth}
Collapse the pointers into one clause-end per \emph{that}, then count.
\begin{center}\footnotesize
\setlength{\tabcolsep}{4pt}\renewcommand{\arraystretch}{1.15}
\begin{longtable}{@{\extracolsep{\fill}}l|ccccccc@{}}
\toprule \textbf{s-op} & the$_{0}$ & boy$_{1}$ & that$_{2}$ & chases$_{3}$ & the$_{4}$ & dog$_{5}$ & swims$_{6}$ \\ \midrule \endfirsthead
\toprule \textbf{s-op} & the$_{0}$ & boy$_{1}$ & that$_{2}$ & chases$_{3}$ & the$_{4}$ & dog$_{5}$ & swims$_{6}$ \\ \midrule \endhead \bottomrule \endfoot
\multicolumn{8}{@{}p{0.98\textwidth}@{}}{\itshape Clause end $e(q)$: $\texttt{rce}$ for a relative \emph{that}, $\texttt{cpe}$ for a sentential one, else $\texttt{idx}$. Here $e(2)=5$: `that chases the dog' closes on \emph{dog}(5).} \\[1pt]
\texttt{e} & 0 & 1 & 5 & 3 & 4 & 5 & 6 \\
\addlinespace[3pt]
\multicolumn{8}{@{}p{0.98\textwidth}@{}}{\itshape $\texttt{where(isthat,\,e,\,-1)}$: mask non-openers to $-1$ so they carry no interval into the count.} \\[1pt]
\texttt{end\_that} & -- & -- & 5 & -- & -- & -- & -- \\
\addlinespace[3pt]
\multicolumn{8}{@{}p{0.98\textwidth}@{}}{\itshape The stabbing count $\#\{\text{\emph{that}}\ q\le i:\ e(q)\ge i\}$ (one attention head). Positions 2--5 are inside the one relative clause $\Rightarrow$ depth 1; the rest 0.} \\[1pt]
\texttt{depth} & 0 & 0 & 1 & 1 & 1 & 1 & 0 \\
\addlinespace[3pt]
\end{longtable}\end{center}
\section{Stage 2 output --- displacement, four flip-typed pieces}
Each piece is a (weighted) stabbing sum over enclosing flips of one type, plus a local head bonus. They add to $\texttt{disp}$.
\begin{center}\footnotesize
\setlength{\tabcolsep}{4pt}\renewcommand{\arraystretch}{1.15}
\begin{longtable}{@{\extracolsep{\fill}}l|ccccccc@{}}
\toprule \textbf{s-op} & the$_{0}$ & boy$_{1}$ & that$_{2}$ & chases$_{3}$ & the$_{4}$ & dog$_{5}$ & swims$_{6}$ \\ \midrule \endfirsthead
\toprule \textbf{s-op} & the$_{0}$ & boy$_{1}$ & that$_{2}$ & chases$_{3}$ & the$_{4}$ & dog$_{5}$ & swims$_{6}$ \\ \midrule \endhead \bottomrule \endfoot
\multicolumn{8}{@{}p{0.98\textwidth}@{}}{\itshape Adverb flip (verb\,$+1$, adverb\,$-1$). All zero --- no adverb.} \\[1pt]
\texttt{d\_adv} & 0 & 0 & 0 & 0 & 0 & 0 & 0 \\
\addlinespace[3pt]
\multicolumn{8}{@{}p{0.98\textwidth}@{}}{\itshape $e-\texttt{idx}+1$: the size of the clause an opener spans (only meaningful at \emph{that}$_2$: $5-2+1=4$).} \\[1pt]
\texttt{span\_clause} & 1 & 1 & 4 & 1 & 1 & 1 & 1 \\
\addlinespace[3pt]
\multicolumn{8}{@{}p{0.98\textwidth}@{}}{\itshape $-\texttt{depth}+(\text{\emph{that}}?\ \texttt{span\_clause}:0)$. Every token in a clause body gets $-1$ per enclosing \emph{that}; the \emph{that} itself gets $+|\text{clause}|$. (Here no \emph{sentential} that, so the clause body $-1$'s come via the relative-clause piece below instead --- $\texttt{d\_cp}$ is 0 because $\texttt{is\_c}$ is empty and depth's sign is carried by $\texttt{d\_nprel}$.)} \\[1pt]
\texttt{d\_cp} & 0 & 0 & 3 & -- & -- & -- & 0 \\
\addlinespace[3pt]
\multicolumn{8}{@{}p{0.98\textwidth}@{}}{\itshape $e[i{+}1]$: the clause end of a \emph{that} sitting at $i{+}1$. Read by a noun to size the relative clause it is about to be modified by.} \\[1pt]
\texttt{e\_at\_p1} & 1 & 5 & 3 & 4 & 5 & 6 & 0 \\
\addlinespace[3pt]
\multicolumn{8}{@{}p{0.98\textwidth}@{}}{\itshape A common noun immediately before a relative \emph{that} gets $+|\text{CP\_rel}|$. \emph{boy}(1): $e[2]-2+1=4$. This is the noun jumping over its relative clause.} \\[1pt]
\texttt{core\_noun\_bonus} & 0 & 4 & 0 & 0 & 0 & 0 & 0 \\
\addlinespace[3pt]
\multicolumn{8}{@{}p{0.98\textwidth}@{}}{\itshape Same bonus for the adjective of an $Adj\,N$ core carrying a relative clause. Zero here (no adjective).} \\[1pt]
\texttt{core\_adj\_bonus} & 0 & 0 & 0 & 0 & 0 & 0 & 0 \\
\addlinespace[3pt]
\multicolumn{8}{@{}p{0.98\textwidth}@{}}{\itshape $\texttt{core\_noun\_bonus}+\texttt{core\_adj\_bonus}$: the $+|\text{RC}|$ for the modified noun. \emph{boy}(1)$=+4$.} \\[1pt]
\texttt{d\_nprel\_core} & 0 & 4 & 0 & 0 & 0 & 0 & 0 \\
\addlinespace[3pt]
\multicolumn{8}{@{}p{0.98\textwidth}@{}}{\itshape $(prev2\_lc=Adj)$: does this relative clause's core noun carry an adjective (so $|\text{core}|=2$)? No.} \\[1pt]
\texttt{has\_adj\_core} & 0 & 0 & 0 & 0 & 0 & 0 & 0 \\
\addlinespace[3pt]
\multicolumn{8}{@{}p{0.98\textwidth}@{}}{\itshape Relative-clause ends with weight-1 core (mask others to $-1$), fed to a stabbing head.} \\[1pt]
\texttt{crel\_end\_w1} & -- & -- & 5 & -- & -- & -- & -- \\
\addlinespace[3pt]
\multicolumn{8}{@{}p{0.98\textwidth}@{}}{\itshape $-(\text{stab}(w1)+2\cdot\text{stab}(w2))$: every token \emph{inside} a relative clause loses $|\text{core}|$. Positions 2--5 each $-1$.} \\[1pt]
\texttt{d\_nprel\_inside} & 0 & 0 & -- & -- & -- & -- & 0 \\
\addlinespace[3pt]
\multicolumn{8}{@{}p{0.98\textwidth}@{}}{\itshape $\texttt{d\_nprel\_core}+\texttt{d\_nprel\_inside}$: the full NP-relative contribution. \emph{boy}$+4$; the clause body $-1$ each.} \\[1pt]
\texttt{d\_nprel} & 0 & 4 & -- & -- & -- & -- & 0 \\
\addlinespace[3pt]
\multicolumn{8}{@{}p{0.98\textwidth}@{}}{\itshape $\texttt{vpe}-\texttt{hb\_end}$: a VP-flip verb's complement size. \emph{chases}(3): $5-3=2$ (`the dog').} \\[1pt]
\texttt{comp\_len} & 0 & 0 & 0 & 2 & 0 & 0 & 0 \\
\addlinespace[3pt]
\multicolumn{8}{@{}p{0.98\textwidth}@{}}{\itshape $(\text{role}=2)\lor(\text{role}=3)$: verbs that actually reverse with a complement. \emph{chases}(3) only.} \\[1pt]
\texttt{is\_vpflip} & 0 & 0 & 0 & 1 & 0 & 0 & 0 \\
\addlinespace[3pt]
\multicolumn{8}{@{}p{0.98\textwidth}@{}}{\itshape VP-flip head (verb, and its adverb) gets $+|\text{comp}|$. \emph{chases}(3)$=+2$.} \\[1pt]
\texttt{d\_vp\_head} & 0 & 0 & 0 & 2 & 0 & 0 & 0 \\
\addlinespace[3pt]
\multicolumn{8}{@{}p{0.98\textwidth}@{}}{\itshape Marks a complement-start: the token whose previous token is a no-adverb VP-flip verb. \emph{the}(4), right after \emph{chases}.} \\[1pt]
\texttt{w1\_here} & 0 & 0 & 0 & 0 & 1 & 0 & 0 \\
\addlinespace[3pt]
\multicolumn{8}{@{}p{0.98\textwidth}@{}}{\itshape That complement's end (from $\texttt{vpe}[i{-}1]$), fed to a stabbing head.} \\[1pt]
\texttt{cs\_end\_w1} & -- & -- & -- & -- & 5 & -- & -- \\
\addlinespace[3pt]
\multicolumn{8}{@{}p{0.98\textwidth}@{}}{\itshape $-(\text{stab}(w1)+2\cdot\text{stab}(w2))$: complement tokens lose $|\text{head}|$ per enclosing VP flip. \emph{the}(4),\emph{dog}(5) each $-1$.} \\[1pt]
\texttt{d\_vp\_comp} & 0 & 0 & 0 & 0 & -- & -- & 0 \\
\addlinespace[3pt]
\multicolumn{8}{@{}p{0.98\textwidth}@{}}{\itshape $\texttt{d\_vp\_head}+\texttt{d\_vp\_comp}$: full transitive/clausal-VP contribution.} \\[1pt]
\texttt{d\_vp} & 0 & 0 & 0 & 2 & -- & -- & 0 \\
\addlinespace[3pt]
\end{longtable}\end{center}
\section{Stage 3 --- assemble and gather}
Sum the pieces to a target slot per token, invert the permutation, gather.
\begin{center}\footnotesize
\setlength{\tabcolsep}{4pt}\renewcommand{\arraystretch}{1.15}
\begin{longtable}{@{\extracolsep{\fill}}l|ccccccc@{}}
\toprule \textbf{s-op} & the$_{0}$ & boy$_{1}$ & that$_{2}$ & chases$_{3}$ & the$_{4}$ & dog$_{5}$ & swims$_{6}$ \\ \midrule \endfirsthead
\toprule \textbf{s-op} & the$_{0}$ & boy$_{1}$ & that$_{2}$ & chases$_{3}$ & the$_{4}$ & dog$_{5}$ & swims$_{6}$ \\ \midrule \endhead \bottomrule \endfoot
\multicolumn{8}{@{}p{0.98\textwidth}@{}}{\itshape $\texttt{d\_adv}+\texttt{d\_cp}+\texttt{d\_nprel}+\texttt{d\_vp}$. This equals the oracle's displacement exactly.} \\[1pt]
\texttt{disp} & 0 & 4 & 2 & 0 & -3 & -3 & 0 \\
\addlinespace[3pt]
\multicolumn{8}{@{}p{0.98\textwidth}@{}}{\itshape $\texttt{idx}+\texttt{disp}$: the HF slot each HI token targets. It is a genuine permutation of $0\dots6$.} \\[1pt]
\texttt{hf\_pos} & 0 & 5 & 4 & 3 & 1 & 2 & 6 \\
\addlinespace[3pt]
\multicolumn{8}{@{}p{0.98\textwidth}@{}}{\itshape $\texttt{kqv(hf\_pos,\,idx,\,idx,\,equals)}$: the \emph{inverse} --- for slot $j$, the source index $i$ with $\texttt{hf\_pos}[i]=j$.} \\[1pt]
\texttt{src} & 0 & 4 & 5 & 3 & 2 & 1 & 6 \\
\addlinespace[3pt]
\multicolumn{8}{@{}p{0.98\textwidth}@{}}{\itshape $\texttt{index\_select(hi\_ids,\,src)}$: gather the token id sitting at each source index into its slot.} \\[1pt]
\texttt{out\_ids} & 0 & 0 & 3 & 33 & 53 & 5 & 45 \\
\addlinespace[3pt]
\end{longtable}\end{center}
Decoding $\texttt{out\_ids}$ back through $\texttt{\_I2W}$ gives the head-final string:
\begin{center}\ttfamily HF = the the dog chases that boy swims\end{center}
\end{document}
